\documentclass[../DoAn.tex]{subfiles}
\begin{document}

Chương này trình bày ba giải pháp nổi bật được hình thành trong quá trình xây dựng WoodCert Auction. Mỗi giải pháp được phân tích theo ba nội dung: bài toán cần giải quyết, quyết định thiết kế và kết quả đạt được. Các chi tiết cài đặt đã trình bày ở Chương~\ref{chapter:Experiment} không được lặp lại; nội dung dưới đây tập trung vào lý do lựa chọn, sự phối hợp giữa các cơ chế và những giới hạn còn tồn tại.

\section{Quản lý trạng thái đấu giá thời gian thực bằng Redis và MySQL}
\label{sec:hybrid-state-management}

\subsection{Giới thiệu bài toán và khó khăn}

Trong phiên đấu giá đang hoạt động, nhiều người tham gia có thể gửi yêu cầu đặt giá trong một khoảng thời gian ngắn. Mỗi yêu cầu phải đồng thời kiểm tra trạng thái phiên, tư cách tham gia, người đang dẫn đầu và bước giá tối thiểu trước khi cập nhật giá mới. Nếu các bước đọc và ghi này bị xen kẽ giữa nhiều yêu cầu, hai lượt đặt giá có thể cùng được đánh giá từ một trạng thái cũ và tạo ra kết quả không nhất quán.

Hệ thống còn phải phản hồi giá mới cho các trình duyệt đang theo dõi mà không yêu cầu người dùng tải lại trang. Ở cuối phiên, hành vi đặt giá trong vài giây cuối có thể làm người tham gia khác không đủ thời gian phản ứng. Do đó, giải pháp cần kết hợp việc kiểm soát cập nhật đồng thời, truyền sự kiện thời gian thực và cơ chế gia hạn phiên. MySQL vẫn phải lưu dữ liệu bền vững để phục vụ truy vấn, đối soát và xử lý sau phiên, trong khi luồng đặt giá không nên phụ thuộc hoàn toàn vào nhiều thao tác khóa liên tiếp trên cơ sở dữ liệu quan hệ.

\subsection{Giải pháp}

WoodCert Auction phân tách vai trò của hai hệ lưu trữ. MySQL quản lý cấu hình phiên, người tham gia, lịch sử đặt giá, bản chụp trạng thái (snapshot) và trạng thái kết thúc. Khi phiên chuyển sang \texttt{ACTIVE}, trạng thái biến đổi nhanh gồm giá hiện tại, người dẫn đầu, thời điểm kết thúc và danh sách người đã đăng ký được khởi tạo trong Redis. Trong giai đoạn này, Redis là nguồn quyết định việc chấp nhận một lượt đặt giá (bid); MySQL tiếp tục là nơi lưu dữ liệu bền vững của hệ thống.

Quyết định này xuất phát từ đặc điểm khác nhau của hai loại dữ liệu. Trạng thái lượt đặt giá đang hoạt động có tần suất đọc ghi cao, vòng đời ngắn và cần một phạm vi thực thi nguyên tử. Ngược lại, thông tin người dùng, sản phẩm, giao dịch và trạng thái kết thúc cần quan hệ, ràng buộc và khả năng truy vấn lâu dài. Nếu chỉ dùng Redis, hệ thống thiếu nguồn dữ liệu quan hệ bền vững để đối soát; nếu chỉ dùng MySQL cho mọi lần cập nhật trạng thái thời gian chạy (runtime), luồng đặt giá phải thực hiện nhiều lần đọc, kiểm tra và ghi trên cùng các bản ghi đang cạnh tranh. Mô hình lai không thay thế vai trò của MySQL mà giới hạn Redis vào phần trạng thái thời gian chạy phù hợp với đặc tính của nó.

Các điều kiện cốt lõi được đóng gói trong một đoạn mã Lua (Lua script) và thực thi nguyên tử trong phạm vi Redis. Đoạn mã lần lượt kiểm tra phiên còn hiệu lực, người dùng thuộc danh sách tham gia, người đặt không phải người đang dẫn đầu và số tiền mới đạt bước giá tối thiểu. Chỉ khi tất cả điều kiện hợp lệ, đoạn mã mới cập nhật giá, người dẫn đầu và mã truy vết của lượt đặt. Cách xử lý này ngăn yêu cầu khác xen vào giữa bước kiểm tra và bước cập nhật trạng thái Redis.

Cơ chế chống đặt giá giây cuối được tích hợp trong cùng thao tác. Nếu thời gian còn lại nhỏ hơn hoặc bằng ngưỡng cấu hình, mặc định 30 giây, đoạn mã cộng thêm thời gian gia hạn, mặc định 60 giây. Đây là cơ chế tạo thêm cơ hội phản hồi cho những người đang tham gia; đồ án không sử dụng cơ chế này để khẳng định chắc chắn rằng giá bán cuối cùng sẽ tăng.

Sau khi Redis chấp nhận lượt đặt giá, phía máy chủ (backend) phát sự kiện \texttt{NEW\_BID} qua WebSocket để các máy khách cập nhật giá và thời gian kết thúc. Lịch sử đặt giá và bản chụp trạng thái phiên sau đó được ghi sang MySQL theo cơ chế cố gắng ghi nhưng không bảo đảm thành công (best-effort) trong cùng luồng xử lý yêu cầu. Lỗi ở bước ghi thứ cấp được lưu vào nhật ký hệ thống nhưng không đảo ngược lượt đặt giá đã được Redis chấp nhận. Quyết định này ưu tiên phản hồi thời gian thực, đồng thời chấp nhận khả năng Redis và MySQL sai lệch tạm thời. Tác vụ nền theo lịch có thể sửa lại bản chụp trạng thái phiên, nhưng không tái tạo được một bản ghi lịch sử đặt giá nếu chính lần ghi bản ghi đó đã thất bại.

Khi phiên đến hạn, tiến trình xử lý vòng đời phiên (lifecycle worker) ưu tiên đọc trạng thái cuối từ Redis. Nếu Redis không khả dụng, hệ thống sử dụng bản chụp trạng thái phiên và lượt đặt giá hợp lệ cao nhất trong MySQL làm đường phục hồi. Kết quả đóng phiên được ghi bền vững trước khi thực hiện quyết toán tiền cọc và tạo đơn hàng. Cách tổ chức này giúp phân biệt rõ trạng thái quyết định lượt đặt giá trong lúc phiên hoạt động với trạng thái nghiệp vụ chính thức sau khi phiên kết thúc. Tuy nhiên, đường phục hồi có thể sử dụng dữ liệu cũ nếu cả bản chụp trạng thái lẫn lịch sử của một lượt đặt giá đều không được ghi thành công; vì vậy đây là cơ chế giảm thiểu sự cố, không phải bảo đảm không mất dữ liệu tuyệt đối.

\subsection{Kết quả đạt được}

Giải pháp đã tạo ra một ranh giới xử lý rõ ràng cho phiên \texttt{ACTIVE}: đoạn mã Lua quyết định và cập nhật lượt đặt giá nguyên tử trong Redis, WebSocket phân phối kết quả, còn MySQL lưu dữ liệu phục vụ vòng đời lâu dài. Kiểm thử tích hợp với Redis thật xác nhận các trường hợp giá thấp, chưa đăng ký, tự đặt đè giá và gia hạn cuối phiên; kiểm thử vòng đời cũng xác nhận tác vụ nền không đóng phiên theo thời điểm cũ sau khi cơ chế chống đặt giá giây cuối đã gia hạn.

Kết quả trên chứng minh tính đúng của các quy tắc đã cài đặt trong phạm vi kiểm thử. Tuy nhiên, đồ án chưa thực hiện đo kiểm tải (benchmark) để đưa ra số liệu về thông lượng hoặc độ trễ, và cơ chế ghi thứ cấp này không tạo ra một giao dịch nguyên tử xuyên Redis và MySQL. Đây là đánh đổi kỹ thuật cần được giám sát bằng nhật ký hệ thống và có thể tiếp tục cải tiến bằng hàng đợi bền vững hoặc cơ chế đối soát lịch sử.

\section{Xử lý ví, tiền cọc và các tác vụ tài chính lũy đẳng}
\label{sec:deposit-settlement-idempotency}

\subsection{Giới thiệu bài toán và khó khăn}

Tiền cọc được dùng để ràng buộc trách nhiệm của người tham gia đấu giá. Khi đăng ký, một phần số dư khả dụng được chuyển sang số dư đóng băng. Sau khi phiên kết thúc, cọc của người không thắng được giải phóng, còn cọc của người thắng được khấu trừ vào giá trị đơn hàng. Các bước tiếp theo còn gồm thanh toán phần còn lại, hoàn tiền, chi trả cho người bán và ghi nhận doanh thu nền tảng.

Một thao tác tài chính có thể bị gọi lại do người dùng gửi lặp yêu cầu, tiến trình thử lại (retry) hoặc tác vụ nền quét lại cùng một đối tượng. Nếu mỗi lần gọi đều làm thay đổi số dư, hệ thống có thể trừ hoặc cộng tiền nhiều lần. Đồng thời, hai nghiệp vụ khác nhau có thể cùng cập nhật một ví trong thời gian gần nhau. Giải pháp vì vậy phải nhận diện thao tác đã xử lý, phát hiện xung đột đồng thời và cho phép tác vụ nền chạy lại mà không lặp biến động tài chính.

\subsection{Giải pháp}

Mỗi biến động ví được gắn với một khóa thao tác nghiệp vụ duy nhất (operation key) được tạo từ loại thao tác và đối tượng liên quan, chẳng hạn phiên đấu giá, người dùng hoặc đơn hàng. Khóa này được lưu trong bảng \texttt{wallet\_operations} và được bảo vệ bằng ràng buộc duy nhất ở cơ sở dữ liệu. Khi một yêu cầu có cùng khóa và cùng nội dung đã hoàn thành thành công, hệ thống bỏ qua việc thay đổi số dư lần nữa. Việc tái sử dụng khóa với dữ liệu khác bị từ chối để tránh một định danh nghiệp vụ đại diện cho hai giao dịch khác nhau.

Khóa được phía máy chủ tạo từ ngữ nghĩa nghiệp vụ thay vì phụ thuộc vào một mã ngẫu nhiên do máy khách cung cấp. Chẳng hạn, thao tác đóng băng cọc của một người dùng trong một phiên luôn ánh xạ tới cùng một khóa. Vì vậy, yêu cầu được gửi lại từ giao diện hoặc lệnh thử lại từ tiến trình nền vẫn được nhận diện là cùng một hành động. Ràng buộc duy nhất trong MySQL đóng vai trò tuyến phòng vệ cuối khi hai yêu cầu đồng thời cùng cố đăng ký một khóa chưa tồn tại.

Bản ghi thao tác theo dõi các trạng thái \texttt{PENDING}, \texttt{SUCCESS} và \texttt{FAILED}. Thao tác chỉ được đánh dấu thành công sau khi giao dịch cơ sở dữ liệu cập nhật ví được cam kết thành công (commit). Nếu giao dịch thất bại, thông tin lỗi được lưu để hỗ trợ truy vết và áp dụng quy tắc thử lại phù hợp. Cơ chế này không phải một sổ cái ký quỹ (escrow) độc lập; nó là nhật ký lũy đẳng phối hợp với số dư khả dụng, số dư đóng băng và các bản ghi giao dịch ví trong MySQL.

Thực thể ví sử dụng thuộc tính phiên bản \texttt{@Version}. Khi hai giao dịch cơ sở dữ liệu cùng đọc một phiên bản và cố ghi các giá trị mới, JPA phát hiện xung đột thay vì để lần ghi sau âm thầm ghi đè lần trước. Khóa thao tác nghiệp vụ giải quyết yêu cầu bị thực hiện lặp, còn khóa lạc quan (optimistic locking) giải quyết các cập nhật cạnh tranh; hai cơ chế bổ sung cho nhau và không thay thế giao dịch của cơ sở dữ liệu.

Ranh giới giao dịch cơ sở dữ liệu cũng được lựa chọn theo phạm vi có thể phục hồi. Một biến động ví và bản ghi giao dịch liên quan được cam kết trong phạm vi nghiệp vụ tương ứng, trong khi việc quyết toán nhiều người tham gia được chia thành các giao dịch \texttt{REQUIRES\_NEW}. Nếu một người tham gia gặp lỗi, các kết quả đã được cam kết của người khác không bị hoàn tác (rollback). Lần chạy sau có thể tiếp tục dựa trên trạng thái người tham gia và khóa thao tác nghiệp vụ. Đánh đổi là hệ thống có thể tạm thời ở trạng thái hoàn tất một phần và cần tác vụ nền tiếp tục xử lý, thay vì đạt tính nguyên tử cho toàn bộ danh sách trong một giao dịch lớn.

Các tác vụ nền sử dụng những cơ chế trên để xử lý vòng đời sau đấu giá. Khi đóng phiên, việc quyết toán cọc của từng người tham gia được thực hiện trong giao dịch riêng để một lỗi cục bộ không làm mất toàn bộ tiến trình. Với đơn quá hạn thanh toán, hệ thống hủy đơn và phân phối khoản cọc bị tịch thu theo cấu hình. Sau khi người mua (Buyer) thanh toán, người bán (Seller) có mặc định 72 giờ để xác nhận đã gửi hoặc bàn giao hàng. Nếu quá hạn xác nhận, tác vụ nền khóa bản ghi, hủy đơn và giao nhận, hoàn lại toàn bộ giá cuối cùng cho người mua rồi đưa sản phẩm về trạng thái có thể tiếp tục xử lý. Mốc 72 giờ này không phải thời hạn kiện hàng phải đến tay người mua.

\subsection{Kết quả đạt được}

Khóa thao tác nghiệp vụ và ràng buộc duy nhất giúp các lệnh tài chính có thể được gọi lại mà không tạo thêm biến động khi thao tác trước đã thành công. Các kiểm thử của dịch vụ ví xác nhận hành vi lũy đẳng, kiểm tra số dư, chuẩn hóa số tiền và phát hiện xung đột phiên bản. Kiểm thử vòng đời thao tác xác nhận việc từ chối khóa có nội dung nghiệp vụ khác, xử lý thao tác đang chờ và cập nhật trạng thái thành công hoặc thất bại.

Các tác vụ nền theo lịch tự động hóa việc quyết toán cọc, xử lý quá hạn thanh toán và bảo vệ người mua khi người bán không xác nhận gửi hàng đúng hạn. Kết quả này giảm phụ thuộc vào thao tác thủ công và cung cấp cơ chế phục hồi có kiểm soát. Dù vậy, việc quyết toán nhiều người tham gia không phải một giao dịch cơ sở dữ liệu duy nhất; hệ thống dựa vào giao dịch tách biệt, trạng thái nghiệp vụ và khóa lũy đẳng để tiếp tục an toàn khi một phần tiến trình gặp lỗi.

\section{Quy trình thẩm định và dấu vân tay tham chiếu của chứng thư}
\label{sec:appraisal-integrity-fingerprint}

\subsection{Giới thiệu bài toán và khó khăn}

Thông tin về loại gỗ, mức định giá và kết luận xác thực ảnh hưởng trực tiếp đến quyết định tham gia đấu giá của người dùng. Vì vậy, người bán không thể tự đưa một sản phẩm chưa qua thẩm định vào đấu giá. Hệ thống cần kiểm soát người được phép tiếp nhận hồ sơ, trạng thái xử lý và thời điểm chứng thư được phát hành.

Bên cạnh việc kiểm soát quy trình, báo cáo cần có một giá trị tham chiếu được tạo từ dữ liệu cốt lõi tại thời điểm phê duyệt. Giá trị này hỗ trợ đối chiếu khi kiểm toán hoặc phát sinh tranh chấp. Tuy nhiên, một mã băm lưu cùng báo cáo không thể tự xác minh chuyên môn của thẩm định viên và cũng không tương đương chữ ký số.

\subsection{Giải pháp}

Sản phẩm bắt đầu ở trạng thái nháp (\texttt{DRAFT}). Sau khi người bán gửi yêu cầu, sản phẩm chuyển sang chờ thẩm định (\texttt{PENDING\_APPRAISAL}); khi thẩm định viên nhận quyền xử lý, trạng thái chuyển thành đang thẩm định (\texttt{UNDER\_APPRAISAL}). Kết quả cuối là \texttt{APPRAISED} hoặc \texttt{REJECTED}; nếu thẩm định viên chủ động trả quyền xử lý, sản phẩm quay lại \texttt{PENDING\_APPRAISAL}. Thao tác nhận hồ sơ khóa bản ghi tại tầng cơ sở dữ liệu và lưu người nhận cùng thời hạn giữ quyền để hạn chế nhiều thẩm định viên xử lý đồng thời. Báo cáo được liên kết với sản phẩm và người thẩm định; chứng thư chỉ được tạo sau khi dữ liệu thẩm định được hoàn tất và báo cáo có định danh ổn định.

Máy trạng thái này phân tách trách nhiệm giữa người bán và thẩm định viên. Người bán cung cấp thông tin sản phẩm và gửi yêu cầu, nhưng không tự chuyển sản phẩm sang trạng thái đã thẩm định. Thẩm định viên tiếp nhận, đánh giá và đưa ra kết luận; các giao diện lập trình ứng dụng (API) tương ứng được giới hạn theo vai trò. Nhờ đó, chứng thư không chỉ là một bản ghi rời rạc mà là kết quả của một chuỗi chuyển trạng thái có điều kiện và có thể truy ngược tới chủ thể xử lý.

Tại thời điểm phê duyệt, hệ thống chuẩn hóa dữ liệu đầu vào (payload) theo thứ tự cố định từ mã sản phẩm, người thẩm định, vật liệu được xác minh, giá trị ước tính, kết luận xác thực, mã chứng thư và thời gian thẩm định. SHA-256 được áp dụng lên dữ liệu này để tạo chuỗi thập lục phân dài 64 ký tự và lưu tại trường \texttt{integrityHash}. Kiểm thử dịch vụ thẩm định tính lại cùng dữ liệu đầu vào và xác nhận giá trị được lưu đúng với kết quả SHA-256 mong đợi.

Việc cố định thứ tự trường và cách biểu diễn là điều kiện cần để cùng một báo cáo luôn sinh ra cùng một dấu vân tay. Nếu định dạng số, thời gian hoặc thứ tự ghép nối thay đổi, giá trị băm cũng thay đổi dù ý nghĩa nghiệp vụ có thể giữ nguyên. Vì vậy, dữ liệu đầu vào hiện tại cần được coi là một quy ước phiên bản của chứng thư; khi mở rộng thêm trường dữ liệu, hệ thống phải duy trì cách tái tạo định dạng cũ hoặc ghi nhận phiên bản thuật toán.

Giải pháp này được xác định là dấu vân tay tham chiếu, không phải cơ chế bảo toàn toàn vẹn độc lập. Phiên bản hiện tại trả mã băm cùng dữ liệu chứng thư nhưng chưa tự động tái tạo dữ liệu đầu vào và so sánh trong luồng truy vấn. Mã băm không xác minh danh tính người ký, không chống chối bỏ và không ngăn người có đủ quyền sửa đồng thời cả báo cáo lẫn giá trị băm.

\subsection{Kết quả đạt được}

Quy trình trạng thái giới hạn việc phát hành chứng thư vào hồ sơ đã được thẩm định viên xử lý, đồng thời lưu mối liên hệ giữa sản phẩm, báo cáo, người thẩm định và thời điểm phê duyệt. Dấu vân tay SHA-256 cung cấp một giá trị tham chiếu có thể dùng cho bước đối chiếu sau này và đã được kiểm tra ở tầng dịch vụ.

Đóng góp chính của giải pháp nằm ở việc kết hợp kiểm soát quy trình thẩm định với dữ liệu chứng thư có thể truy vết. Mức bảo đảm hiện tại vẫn phụ thuộc vào quyền truy cập cơ sở dữ liệu và chất lượng thẩm định ngoài phần mềm. Để nâng thành cơ chế xác thực mạnh hơn, hệ thống cần bổ sung bước kiểm tra mã băm tự động và lưu chữ ký số hoặc bằng chứng xác thực ở một miền tin cậy độc lập.

\section{Tổng kết chương}
\label{sec:conclusion-ch5}

Ba giải pháp trên tập trung vào các điểm khó nhất của WoodCert Auction: cập nhật lượt đặt giá thời gian thực, xử lý tài chính có thể chạy lại an toàn và kiểm soát quy trình thẩm định. Mỗi giải pháp đều xác định rõ phạm vi bảo đảm và đánh đổi còn tồn tại; các kết quả kiểm thử chứng minh hành vi đã cài đặt nhưng không thay thế đo kiểm tải, kiểm toán tài chính hoặc chứng nhận an toàn độc lập.

\end{document}
